\chapter{Sprint 1 : Préparation des données et amorçage}
\label{chap:sprint1}

\sectionsansnum{Introduction}

Ce chapitre rend compte du premier sprint, consacré à la constitution du corpus
d'apprentissage. Aucun modèle ne vaut mieux que les données qui l'alimentent : avant
toute architecture, il fallait transformer des enregistrements bruts hétérogènes en un
corpus homogène et étiqueté. Le sprint s'est heurté à une contrainte externe --- le
délai d'instruction de la demande d'accès aux cohortes cliniques --- que nous avons
choisi d'occuper plutôt que de subir.

\section{Objectif du sprint}

L'objectif était double. Il s'agissait d'abord de \textbf{construire la chaîne de
préparation du signal} : lire un fichier au format EDF quel que soit son constructeur,
en extraire les cinq dérivations utiles, uniformiser la fréquence d'échantillonnage,
découper le signal en fenêtres conformes au référentiel AASM et neutraliser les écarts
d'amplitude entre sujets.

Il s'agissait ensuite de \textbf{valider cette chaîne de bout en bout} sans attendre
l'autorisation d'accès aux cohortes du \emph{National Sleep Research
Resource}~\cite{zhang2018nsrr}, dont l'instruction prend plusieurs semaines. Le jeu
public Sleep-EDF~\cite{kemp2000}, diffusé par PhysioNet~\cite{goldberger2000} et
immédiatement téléchargeable, a servi de banc d'essai.

Ce choix mérite d'être explicité, car il constitue le principal arbitrage de gestion de
projet du stage. Attendre l'autorisation aurait immobilisé quatre semaines. Travailler
sur un jeu public de substitution permettait de valider la lecture des fichiers, la
résolution des dérivations, la segmentation et la boucle d'entraînement avant même de
disposer des données définitives. Lorsque l'autorisation est arrivée en semaine 5, la
chaîne était éprouvée et il n'a plus été question que de la rejouer à plus grande
échelle.

\section{Analyse du sprint 1}

\subsection{Sprint Backlog}

Le sprint a embarqué quatre récits utilisateur pour un total de 29 points de
complexité, détaillés au tableau~\ref{tab:bl-s1}.

\begin{table}[H]
\centering\small
\renewcommand{\arraystretch}{1.15}
\begin{tabular}{@{}c|L{2.5cm}|L{4.6cm}|L{4.6cm}|c@{}}
\hline
\ligneentete \entete{ID} & \entete{Module} & \entete{R\'ecit utilisateur} & \entete{T\^aches} & \entete{Estim.} \\
\hline
\multirow{3}{*}{US-22} & \multirow{3}{=}{Lecture des enregistrements} & \multirow{3}{=}{En tant que syst\`eme, je veux lire un fichier EDF et en extraire les d\'erivations physiologiques} & Int\'egrer la biblioth\`eque de lecture des formats physiologiques & \multirow{3}{*}{8 pts} \\
 &  &  & Inventorier les d\'erivations et la fr\'equence native &  \\
 &  &  & G\'erer les fichiers tronqu\'es ou illisibles &  \\
\hline
\multirow{3}{*}{US-23} & \multirow{3}{=}{R\'esolution des canaux} & \multirow{3}{=}{En tant que syst\`eme, je veux r\'esoudre les canaux malgr\'e l'h\'et\'erog\'en\'eit\'e des conventions de nommage} & Construire la table d'alias des cinq emplacements logiques & \multirow{3}{*}{5 pts} \\
 &  &  & D\'efinir la strat\'egie de repli sur l'EEG principal &  \\
 &  &  & Exiger le montage complet \`a l'entra\^inement &  \\
\hline
\multirow{4}{*}{US-24} & \multirow{4}{=}{Normalisation du signal} & \multirow{4}{=}{En tant que syst\`eme, je veux r\'e\'echantillonner, segmenter en \'epoques de 30 s et normaliser} & R\'e\'echantillonner \`a 100\,Hz par m\'ethode de Fourier & \multirow{4}{*}{8 pts} \\
 &  &  & D\'ecouper en fen\^etres de 3\,000 \'echantillons &  \\
 &  &  & Centrer et r\'eduire chaque fen\^etre &  \\
 &  &  & S\'erialiser en tenseurs par sujet &  \\
\hline
\multirow{3}{*}{US-25} & \multirow{3}{=}{Amor\c{c}age} & \multirow{3}{=}{En tant qu'ing\'enieur, je veux valider la cha\^ine sur un jeu public avant l'acc\`es aux cohortes cliniques} & T\'el\'echarger et pr\'eparer le jeu public & \multirow{3}{*}{8 pts} \\
 &  &  & Entra\^iner un mod\`ele de faisabilit\'e &  \\
 &  &  & V\'erifier la boucle compl\`ete de bout en bout &  \\
\hline
\end{tabular}
\caption{Sprint Backlog --- Sprint 1}
\label{tab:bl-s1}
\end{table}

\subsection{Diagramme des cas d'utilisation}

Ce sprint ne produit pas de fonctionnalité directement visible par le praticien : il
construit le socle sur lequel reposeront tous les suivants. Le cas d'utilisation
principal est donc porté par le moteur d'inférence, acteur secondaire du système
(figure~\ref{fig:uc-s1}).

\begin{figure}[H]
\centering
\begin{tikzpicture}[
    x=1cm,y=1cm,
    uc/.style     = {ellipse, draw=hypnoblue!60, fill=hypnolight,
                     align=center, font=\scriptsize, text width=2.6cm,
                     inner sep=1.5pt, minimum height=1.0cm},
    ucp/.style    = {ellipse, draw=hypnored!70, fill=hypnored!7,
                     align=center, font=\scriptsize\bfseries, text width=2.6cm,
                     inner sep=1.5pt, minimum height=1.0cm},
    lien/.style   = {hypnogray!75, line width=0.55pt},
    rel/.style    = {-{Latex[length=2mm]}, hypnoblue!70, line width=0.5pt, dashed},
    et/.style     = {font=\tiny\itshape, text=hypnoblue!85, fill=white, inner sep=1pt},
    nomact/.style = {font=\scriptsize\bfseries, align=center, text width=2.2cm},
  ]

  \tikzset{acteur/.pic={
      \draw[hypnoblue,line width=0.8pt] (0,0) circle (0.17);
      \draw[hypnoblue,line width=0.8pt] (0,-0.17) -- (0,-0.66);
      \draw[hypnoblue,line width=0.8pt] (-0.28,-0.34) -- (0.28,-0.34);
      \draw[hypnoblue,line width=0.8pt] (0,-0.66) -- (-0.24,-1.04);
      \draw[hypnoblue,line width=0.8pt] (0,-0.66) -- (0.24,-1.04);}}

  % Frontiere du systeme
  \draw[hypnoblue!55,line width=0.9pt,rounded corners=3pt]
       (3.1,-6.0) rectangle (12.5,2.4);
  \node[font=\small\bfseries,text=hypnoblue,anchor=north west] at (3.35,2.25)
       {Module de préparation des données};

  % Cas principal et cas terminal, colonne de gauche
  \node[ucp] (prep) at (5.1, 0.9)  {Préparer un\\enregistrement};
  \node[uc]  (ser)  at (5.1,-4.3)  {Sérialiser le\\corpus};

  % Sous-cas inclus, colonne de droite
  \node[uc]  (lire) at (10.0, 1.2) {Décoder le\\fichier EDF};
  \node[uc]  (res)  at (10.0,-0.3) {Résoudre les\\dérivations};
  \node[uc]  (ree)  at (10.0,-1.8) {Rééchantillonner\\à 100\,Hz};
  \node[uc]  (seg)  at (10.0,-3.3) {Segmenter en\\époques de 30\,s};
  \node[uc]  (nor)  at (10.0,-4.8) {Normaliser\\chaque époque};

  % Acteur
  \pic at (1.3,0.6) {acteur};
  \node[nomact,anchor=north] at (1.3,-0.05) {Moteur\\d'inférence};

  \draw[lien] (1.7,0.25) -- (prep.west);
  \draw[lien] (1.7,0.25) -- (ser.west);

  \draw[rel] (prep) -- node[et,pos=0.62] {$\ll$include$\gg$} (lire);
  \draw[rel] (prep) -- node[et,pos=0.62] {$\ll$include$\gg$} (res);
  \draw[rel] (prep) -- node[et,pos=0.62] {$\ll$include$\gg$} (ree);
  \draw[rel] (prep) -- node[et,pos=0.66] {$\ll$include$\gg$} (seg);
  \draw[rel] (prep) -- node[et,pos=0.70] {$\ll$include$\gg$} (nor);

\end{tikzpicture}
\caption{Diagramme des cas d'utilisation du sprint 1 --- préparation des données}
\label{fig:uc-s1}
\end{figure}

\subsubsection{Description textuelle des cas d'utilisation}

\begin{casutilisation}{Description textuelle du cas « Préparer un enregistrement »}
Titre & Préparer un enregistrement \\ \hline
Acteur & Moteur d'inférence \\ \hline
Description & Transformer un enregistrement polysomnographique brut en un ensemble de
fenêtres normalisées, directement exploitables par un modèle de classification \\ \hline
Préconditions & Un fichier au format EDF est disponible et lisible \\ \hline
Scénario de base &
\begin{minipage}[t]{\linewidth}
\begin{enumerate}[leftmargin=1.1em,itemsep=1pt,topsep=3pt]
  \item Le système décode le fichier et inventorie ses dérivations.
  \item Il apparie chaque dérivation attendue à son équivalent dans le fichier.
  \item Il ramène l'ensemble des signaux à \SI{100}{\hertz}.
  \item Il découpe le signal en fenêtres de \num{3000} échantillons.
  \item Il centre et réduit chaque fenêtre indépendamment.
  \item Il sérialise le résultat sous forme de tenseur accompagné de ses étiquettes.
\end{enumerate}
\end{minipage} \\ \hline
Scénario alternatif &
\begin{minipage}[t]{\linewidth}
\begin{itemize}[leftmargin=1.1em,itemsep=1pt,topsep=3pt]
  \item \textbf{2a} --- Une dérivation est absente : en production, repli documenté sur
        l'EEG principal ; à l'entraînement, le sujet est écarté du corpus.
  \item \textbf{4a} --- L'enregistrement est plus court qu'une époque : le traitement
        s'interrompt avec un message explicite.
  \item \textbf{6a} --- L'hypnogramme de référence est absent : le sujet est écarté.
\end{itemize}
\end{minipage}
\end{casutilisation}

\subsection{Diagramme de séquence objet}

La figure~\ref{fig:seq-s1} détaille l'enchaînement des traitements appliqués à un
enregistrement.

\begin{figure}[H]
\centering
\resizebox{0.96\textwidth}{!}{%
\begin{sequencediagram}
  \newthread[hypnolight]{p}{\scriptsize Script de préparation}
  \newinst[2.0]{l}{\scriptsize Lecteur EDF}
  \newinst[2.0]{a}{\scriptsize Résolveur de canaux}
  \newinst[2.0]{s}{\scriptsize Traitement du signal}
  \newinst[2.0]{x}{\scriptsize Analyseur XML}
  \newinst[2.0]{f}{\scriptsize Corpus sérialisé}

  \begin{call}{p}{\scriptsize ouvrir l'enregistrement}{l}{\scriptsize dérivations, fréquence}
  \end{call}
  \begin{call}{p}{\scriptsize apparier les dérivations}{a}{\scriptsize table de correspondance}
  \end{call}
  \begin{sdblock}{\scriptsize Contrôle}{\scriptsize montage complet exigé}
    \begin{call}{p}{\scriptsize écarter le sujet si incomplet}{p}{}
    \end{call}
  \end{sdblock}
  \begin{call}{p}{\scriptsize rééchantillonner}{s}{\scriptsize signal à 100\,Hz}
  \end{call}
  \begin{call}{p}{\scriptsize lire l'hypnogramme}{x}{\scriptsize séquence de stades}
  \end{call}
  \begin{call}{p}{\scriptsize segmenter et normaliser}{s}{\scriptsize tenseur d'époques}
  \end{call}
  \begin{call}{p}{\scriptsize sérialiser}{f}{\scriptsize fichier par sujet}
  \end{call}
\end{sequencediagram}}
\caption{Diagramme de séquence objet du sprint 1 --- préparation d'un enregistrement}
\label{fig:seq-s1}
\end{figure}

\section{Réalisation}

\subsection{La chaîne de préparation}

La chaîne réalisée comporte six étapes, chacune répondant à une source de variabilité
identifiée (figure~\ref{fig:chaine}).

\begin{figure}[H]
\centering
\begin{tikzpicture}[
    x=1cm,y=1cm,
    et/.style  = {rectangle,rounded corners=2pt,draw=hypnoblue!60,fill=hypnolight,
                  align=center,font=\scriptsize,minimum height=1.25cm,
                  text width=2.15cm},
    fl/.style  = {-{Latex[length=2mm]},hypnoblue!75,line width=0.8pt},
    note/.style= {font=\tiny\itshape,text=hypnogray,align=center,text width=2.3cm},
  ]

  \node[et] (a) at (0.0,0)  {Décodage\\du fichier};
  \node[et] (b) at (2.75,0) {Résolution\\des dérivations};
  \node[et] (c) at (5.5,0)  {Rééchantillonnage\\à 100\,Hz};
  \node[et] (d) at (8.25,0) {Segmentation\\en époques};
  \node[et] (e) at (11.0,0) {Normalisation\\par époque};
  \node[et] (f) at (13.75,0){Tenseur\\$(N, C, 3000)$};

  \foreach \x/\y in {a/b,b/c,c/d,d/e,e/f} \draw[fl] (\x) -- (\y);

  \node[note,anchor=north] at (0.0,-0.78)  {format\\constructeur};
  \node[note,anchor=north] at (2.75,-0.78) {conventions de\\nommage};
  \node[note,anchor=north] at (5.5,-0.78)  {fréquences\\d'acquisition};
  \node[note,anchor=north] at (8.25,-0.78) {convention\\AASM 30\,s};
  \node[note,anchor=north] at (11.0,-0.78) {amplitudes\\inter-sujets};
  \node[note,anchor=north] at (13.75,-0.78){entrée des\\modèles};

  \node[font=\tiny\bfseries,text=hypnored,anchor=east] at (-1.05,-1.35)
       {variabilité traitée :};

\end{tikzpicture}
\vspace{0.3cm}
\caption[Chaîne de traitement du signal]%
        {Chaîne de préparation : chaque étape neutralise une source de variabilité.}
\label{fig:chaine}
\end{figure}

Quatre décisions de traitement méritent d'être explicitées.

\paragraph{Sélection stricte des dérivations.} À l'entraînement, les cinq dérivations
sont exigées sous leur dénomination exacte et tout écart écarte le sujet. En production,
au contraire, une table d'alias couvre les conventions de nommage des différents
constructeurs et replie sur l'EEG principal une dérivation absente. Cette asymétrie est
voulue : sélectivité pour la pureté du corpus d'apprentissage, tolérance pour
l'interopérabilité clinique. C'est un des rares endroits du système où une même fonction
est délibérément implémentée deux fois différemment.

\paragraph{Fusion des anciens stades 3 et 4.} Les annotations des cohortes suivent
encore la nomenclature de Rechtschaffen et Kales~\cite{rechtschaffen1968}. Le stade 4
est fusionné dans le N3, conformément à la révision AASM de 2007.

\paragraph{Comblement des époques non scorées.} Les fenêtres sans stade attribué
reçoivent le dernier stade valide observé ; celles qui précèdent le premier stade scoré
sont considérées comme de l'éveil. Cette dernière convention est une hypothèse ---
l'enregistrement débute avant l'extinction des lumières --- et non une donnée.

\paragraph{Normalisation locale.} La statistique de centrage et de réduction est
calculée sur chaque fenêtre isolément, et non sur le corpus. Ce choix neutralise les
écarts d'amplitude liés à l'impédance des électrodes, au gain de l'amplificateur et à la
morphologie du sujet, ainsi que la dérive lente au cours de la nuit. Il a surtout une
conséquence méthodologique : \textbf{aucune statistique globale n'étant estimée, aucune
information ne peut fuir de l'ensemble d'entraînement vers l'ensemble de test par la
normalisation}. En contrepartie, l'information d'amplitude absolue est perdue, alors
qu'elle caractérise partiellement le sommeil profond.

\subsection{Le corpus produit}

La préparation est assurée par un script exécuté une seule fois, dont le produit
alimentera les douze entraînements du sprint suivant. Chaque enregistrement est apparié
à son fichier d'annotations, ses cinq dérivations sont extraites, le signal est ramené à
\SI{100}{\hertz}, découpé en fenêtres de \num{3000} échantillons et centré-réduit fenêtre
par fenêtre. Le résultat est sérialisé sous forme d'un tenseur en simple précision
accompagné d'un vecteur d'étiquettes en entier court, soit environ \SI{60}{\mega\octet}
par nuit.

Le script est \textbf{relançable} : un sujet déjà converti est ignoré, ce qui permet
d'étaler la campagne de préparation sur plusieurs sessions. Détail d'ingénierie modeste,
mais c'est exactement ce qui rend une expérience reproductible.

Cent quatre-vingt-seize enregistrements de la cohorte SHHS1~\cite{quan1997} ont été
retenus, pour un total de \num{195130} époques de 30 secondes. Il faut souligner que le
corpus est \textbf{filtré et non échantillonné} : un sujet est écarté s'il lui manque une
des cinq dérivations requises ou son fichier d'annotations. Les sujets retenus sont donc
les premiers de la cohorte à satisfaire le montage complet, ce qui introduit un biais de
sélection dont nous discuterons la portée.

\begin{figure}[H]
\centering
\begin{tikzpicture}[x=1cm,y=1cm]

  \def\W{11.4}

  % ── Barre 5 classes ───────────────────────────────────────────────────────
  \node[font=\scriptsize\bfseries,text=hypnogray,anchor=east] at (-0.25,1.55)
       {5 classes};
  \foreach \a/\b/\col/\lab/\pct in {
      0/30.70/stWake/W/30{,}7,
      30.70/33.96/stN1/N1/3{,}3,
      33.96/74.13/stN2/N2/40{,}2,
      74.13/87.09/stN3/N3/13{,}0,
      87.09/100/stREM/REM/12{,}9}
  {
    \pgfmathsetmacro{\xa}{\a*\W/100}
    \pgfmathsetmacro{\xb}{\b*\W/100}
    \pgfmathsetmacro{\xm}{(\xa+\xb)/2}
    \fill[\col!75] (\xa,1.15) rectangle (\xb,1.95);
    \draw[white,line width=0.8pt] (\xb,1.15) -- (\xb,1.95);
    \node[font=\tiny\bfseries,text=white] at (\xm,1.68) {\lab};
    \node[font=\tiny,text=white] at (\xm,1.40) {\pct\,\%};
  }

  % ── Barre 3 classes ───────────────────────────────────────────────────────
  \node[font=\scriptsize\bfseries,text=hypnogray,anchor=east] at (-0.25,0.35)
       {3 classes};
  \foreach \a/\b/\col/\lab/\pct in {
      0/30.70/stWake/Wake/30{,}7,
      30.70/87.09/stN2/NREM/56{,}4,
      87.09/100/stREM/REM/12{,}9}
  {
    \pgfmathsetmacro{\xa}{\a*\W/100}
    \pgfmathsetmacro{\xb}{\b*\W/100}
    \pgfmathsetmacro{\xm}{(\xa+\xb)/2}
    \fill[\col!75] (\xa,-0.05) rectangle (\xb,0.75);
    \draw[white,line width=0.8pt] (\xb,-0.05) -- (\xb,0.75);
    \node[font=\tiny\bfseries,text=white] at (\xm,0.48) {\lab};
    \node[font=\tiny,text=white] at (\xm,0.20) {\pct\,\%};
  }

  % ── Annotation ────────────────────────────────────────────────────────────
  \draw[hypnored,line width=0.7pt] (3.50,2.10) -- (3.50,2.45) -- (4.30,2.45);
  \node[font=\scriptsize,text=hypnored,anchor=west] at (4.35,2.45)
       {le N1 ne représente que \num{6361} époques sur \num{195130}};

  \node[font=\tiny\itshape,text=hypnogray,anchor=north] at (5.7,-0.35)
       {corpus SHHS1 --- 196 enregistrements, \num{195130} époques de 30 secondes};

\end{tikzpicture}
\vspace{0.35cm}
\caption[Distribution des stades dans le corpus SHHS1]%
        {Distribution des stades dans le corpus d'entraînement, aux deux granularités.}
\label{fig:distribution}
\end{figure}

Le déséquilibre est frappant et conditionne tout le sprint suivant. Le N2 occupe à lui
seul \SI{40,2}{\percent} des époques, tandis que le \textbf{N1 n'en représente que
\SI{3,3}{\percent}}, soit \num{6361} exemples. Un modèle qui ignorerait purement et
simplement cette classe perdrait à peine trois points d'exactitude globale : c'est
exactement le comportement que la pondération des classes devra empêcher.

\subsection{Découpage et pondération}

Le corpus est découpé en trois parts stratifiées sur la classe, dans les proportions
\num{80} / \num{10} / \num{10}, soit \num{156104} époques d'entraînement et \num{19513}
pour chacun des ensembles de validation et de test. La pondération des classes est fixée
à l'inverse de leur fréquence, puis renormalisée : le REM se voit attribuer un poids
\num{4,4} fois supérieur à celui du NREM en trois classes, et le N1 un poids \num{12,3}
fois supérieur à celui du N2 en cinq classes.

\sectionsansnum{Conclusion}

Ce premier sprint a produit le socle de tout le projet : une chaîne de préparation
opérationnelle et un corpus de \num{195130} époques étiquetées, issues de 196
enregistrements cliniques.

L'arbitrage de départ --- travailler sur un jeu public pendant l'instruction de la
demande d'accès --- s'est révélé payant : lorsque les données définitives sont arrivées,
la chaîne était éprouvée et il n'a suffi que de la rejouer à plus grande échelle.

Deux constats issus de ce sprint pèseront sur les suivants : le corpus est fortement
déséquilibré, avec un stade N1 marginal, et il résulte d'un filtrage sur la complétude
du montage plutôt que d'un tirage aléatoire. Le sprint suivant exploite ce corpus pour
entraîner les trois familles d'architectures de classification.
